% Nêu tổng quan phương pháp thực hiện

\begin{enumerate}
    \item \textbf{Thu thập và phân tích dữ liệu:}
    
    Thu thập các bộ dữ liệu lưu lượng mạng từ các nguồn công khai (như Abilene, GÉANT) hoặc \textit{(for Kiet, if you dont want this then remove)} chủ động tạo tập dữ liệu tổng hợp (Synthetic Data) bằng mô hình hấp dẫn (Gravity Model) cho topo NSFNET. Tập dữ liệu bao gồm thông tin về băng thông, độ trễ, tỷ lệ mất gói và mức sử dụng liên kết. Tiến hành tiền xử lý dữ liệu bằng cách loại bỏ giá trị bất thường, chuẩn hóa dữ liệu, tạo cửa sổ thời gian (Sliding window) và trích xuất các đặc trưng cần thiết phục vụ cho quá trình huấn luyện mô hình. Dữ liệu được phân chia theo 3 miền mạng để mô phỏng quá trình huấn luyện phân tán trong Federated Learning.

    \item \textbf{Môi trường triển khai:}

    Hệ thống được triển khai trên nền tảng Ubuntu với ngôn ngữ Python. Quá trình triển khai chia làm hai giai đoạn: \textit{(For me/Kiet)} Giai đoạn huấn luyện sử dụng môi trường mô phỏng Gym (OpenAI Gymnasium) tùy chỉnh với mô hình tính toán độ trễ M/M/1 để tăng tốc độ huấn luyện RL; Giai đoạn đánh giá và demo sử dụng Mininet với topo NSFNET và bộ điều khiển Ryu (OpenFlow 1.3). Mô hình dự đoán lưu lượng dùng PyTorch, huấn luyện Federated Learning với framework Flower (FedAvg). Thuật toán Reinforcement Learning (PPO) được triển khai để tối ưu quyết định định tuyến dựa trên trạng thái mạng và kết quả dự đoán lưu lượng, nâng cao hiệu quả mạng ISP.

    \item \textbf{Xây dựng mô hình dự đoán lưu lượng:}
    
    Thiết kế mô hình dự đoán lưu lượng mạng sử dụng mạng nơ-ron hồi tiếp LSTM (Long Short-Term Memory) nhằm dự đoán lưu lượng trong các khoảng thời gian tiếp theo dựa trên dữ liệu lịch sử. Đầu ra của mô hình là giá trị lưu lượng dự đoán trong các khoảng thời gian tiếp theo trên các liên kết mạng. Kết quả dự đoán này được sử dụng làm thông tin đầu vào cho mô-đun Traffic Engineering, kết hợp với trạng thái mạng hiện tại để hỗ trợ tác tử RL đưa ra quyết định điều khiển lưu lượng chủ động.

    \item \textbf{Xây dựng cơ chế tối ưu Traffic Engineering:}
    
    Xây dựng cơ chế Traffic Engineering sử dụng Reinforcement Learning (PPO). \textit{(For me)} Bài toán được mô hình hóa thành MDP: Trạng thái (State) gồm 104 chiều (mức sử dụng 42 liên kết hiện tại, 42 dự đoán tương lai, 20 nhu cầu lưu lượng); Hành động (Action) là không gian rời rạc (MultiDiscrete) chọn 1 trong 3 đường đi (K=3) cho top 20 luồng giao thông lớn nhất; Phần thưởng (Reward) được tính bằng hàm phạt trọng số dựa trên MLU (Maximum Link Utilization), độ trễ đầu cuối và tỷ lệ mất gói (dựa trên mô hình hàng đợi M/M/1). Quyết định điều khiển sẽ được chuyển đến SDN Controller để thực thi thông qua việc cập nhật các quy tắc chuyển tiếp (Flow Rules) trên mạng.

    \item \textbf{Tích hợp và triển khai hệ thống:}
    
    Tích hợp mô hình Traffic Prediction dựa trên Federated Learning với mô-đun Traffic Engineering sử dụng Reinforcement Learning trong kiến trúc SDN. Kết quả dự đoán lưu lượng được kết hợp với trạng thái mạng hiện tại để hỗ trợ RL đưa ra quyết định điều khiển lưu lượng. Các quyết định này được chuyển đến SDN Controller để thực thi trên topo NSFNET trong môi trường Mininet.

    \item \textbf{Thực nghiệm và đánh giá:}
    
    Thực hiện các kịch bản thực nghiệm trên topo NSFNET với nhiều mức tải mạng khác nhau, bao gồm lưu lượng bình thường, lưu lượng tăng đột biến và các trường hợp xảy ra tắc nghẽn trên một hoặc nhiều liên kết. Hệ thống được đánh giá thông qua các chỉ số dự đoán lưu lượng như MAE, RMSE, MAPE và các chỉ số Traffic Engineering như độ trễ đầu cuối, thông lượng, tỷ lệ sử dụng liên kết, tỷ lệ mất gói, mức độ cân bằng tải, thời gian hội tụ và các thông số khác. Kết quả được so sánh với các phương pháp định tuyến truyền thống và mô hình học tập trung để đánh giá khả năng cải thiện độ chính xác dự đoán và hiệu quả điều khiển lưu lượng.
\end{enumerate}